% Regression: t2_rgfp4 — imported current-behavior fixture from the handoff corpus.
% Formula: e (U a unitary on the four physical indices of a 2x2 plaquette,
% t2_rgfp4 — RMP arXiv:2011.12127, caption FigPEPSRG, file
% fig2_fig2-pepe_Figboundary-theoryb-1 as it now stands: the four
% renormalization-fixed-point possibilities for a 2D PEPS tensor.
%
% Formulae (U a unitary on the four physical indices of a 2x2 plaquette,
% invertibles on the virtual ones):
%  (a) the object renormalized: the 2x2 plaquette of PEPS tensors;
%  (b) U (plaquette) = A (x) |phi_1>|phi_2>|phi_3> — the original tensor
%      tensor-multiplied by DISENTANGLED single-site product states, each
%      keeping one dangling leg;
%  (c) U (plaquette) = A (x) M — the original tensor times a leftover
%      matrix-product (two bonded 3-leg tensors);
%  (d) U (plaquette) = A (x) A (x) tr-bead — two copies of the original
%      tensor, times a scalar drawn as a bead whose virtual index closes on
%      itself (the RMP's bracketed bead; here the hooks trace closure).
%
% Declarative reading: the plaquette and every copy of A are 1x1 / 2x2
% oblique lattice windows (boundary legs = the plaquette's virtual legs,
% physical=up); the product-state and matrix factors are one-row grids
% sealed per side (west=none / east=none); the traced bead is
% periodic, trace style=hooks on a single cell.  The (x) symbols are algebra
% between pictures, not ink.
%
% Noted gaps (cosmetic): the RMP slants the disentangled product beads'
% single legs at oblique angles; the grid offers up / west / east
% directions only, so the three beads in (b) coarsen to those.  The RMP
% hue distinctions are absent (monochrome ink).
\documentclass[varwidth=19cm, border=6pt]{standalone}
\usepackage{amsmath, amssymb}
\usepackage{tenkz}
\begin{document}

(a) the renormalized plaquette:
\[
  \begin{tenkzlattice}[rows=2, cols=2, frame=oblique, boundary legs,
                       physical=up]
  \end{tenkzlattice}
\]

(b) original tensor $\otimes$ disentangled product states:
\[
  \begin{tenkzlattice}[rows=1, cols=1, frame=oblique, boundary legs,
                       physical=up]
  \end{tenkzlattice}
  \;\otimes\;
  % three 1-leg beads: leg up / leg east / leg west
  \begin{tenkz}[physical=up, west=none, east=none]
    \tn{}
  \end{tenkz}
  \;
  \begin{tenkz}[west=none]
    \tn{}
  \end{tenkz}
  \;
  \begin{tenkz}[east=none]
    \tn{}
  \end{tenkz}
\]

(c) original tensor $\otimes$ a generating matrix product:
\[
  \begin{tenkzlattice}[rows=1, cols=1, frame=oblique, boundary legs,
                       physical=up]
  \end{tenkzlattice}
  \;\otimes\;
  % two bonded 3-leg tensors, outer virtual sides sealed
  \begin{tenkz}[physical=up, west=none, east=none]
    \tn{} & \tn{}
  \end{tenkz}
\]

(d) two copies of the original tensor $\otimes$ the traced bead:
\[
  \begin{tenkzlattice}[rows=1, cols=1, frame=oblique, boundary legs,
                       physical=up]
  \end{tenkzlattice}
  \;\otimes\;
  \begin{tenkzlattice}[rows=1, cols=1, frame=oblique, boundary legs,
                       physical=up]
  \end{tenkzlattice}
  \;\otimes\;
  % the bead whose virtual wire closes on itself: hooks trace closure
  \begin{tenkz}[physical=up, periodic, trace style=hooks]
    \tn{}
  \end{tenkz}
\]
\end{document}
